\documentclass[11pt,a4paper]{article}
\usepackage[utf8]{inputenc}
\usepackage[T1]{fontenc}
\usepackage{amsmath,amssymb}
\usepackage{algorithm}
\usepackage{algpseudocode}
\usepackage{listings}
\usepackage{xcolor}
\usepackage{hyperref}
\usepackage{graphicx}
\usepackage{booktabs}
\usepackage{caption}
\usepackage{geometry}
\geometry{margin=1in}

\title{\textbf{UMI: Universal Mahjong Interface Protocol}\\
\large A Communication Protocol for Japanese Riichi Mahjong Engines\\
Version 1.0 Specification}
\author{
    \textbf{Protocol Design Committee}\\
    Research Institute of Computer Games\\
    \texttt{umi-spec@riichi.org}
}
\date{March 2025}

\lstset{
    basicstyle=\ttfamily\small,
    keywordstyle=\color{blue},
    commentstyle=\color{gray},
    stringstyle=\color{red},
    frame=single,
    breaklines=true,
    showstringspaces=false
}

\begin{document}

\maketitle

\begin{abstract}
This document specifies the Universal Mahjong Interface (UMI) protocol, a text-based communication standard for Japanese Riichi Mahjong engines and user interfaces. Inspired by the Universal Chess Interface (UCI), UMI enables interoperability between different Mahjong AI engines and graphical or command-line interfaces. The protocol supports local gameplay via named pipes, networked gameplay via TCP sockets, and decentralized peer-to-peer discovery via WebSocket over Distributed Hash Table (DHT) networks. We present the complete command set, state representation, move notation, and network architecture for building distributed Mahjong playing systems.
\end{abstract}

\section{Introduction}

Japanese Riichi Mahjong is a complex four-player game of incomplete information, featuring intricate scoring rules (yaku), tile interactions (calls), and strategic depth. Unlike chess, which is two-player and perfect-information, Mahjong presents unique challenges for protocol design: simultaneous actions, hidden information, multiple win conditions, and four-way interaction.

The Universal Mahjong Interface (UMI) addresses these challenges by providing:
\begin{itemize}
    \item \textbf{Stateless engine design}: Engines receive complete game state on each decision
    \item \textbf{Multi-transport support}: Named pipes, TCP sockets, and WebSocket/DHT
    \item \textbf{Decentralized matchmaking}: Peer discovery without central servers
    \item \textbf{Complete information hiding}: Protocol respects Mahjong's hidden hand nature
\end{itemize}

\section{Protocol Architecture}

\subsection{Design Philosophy}

Following UCI's successful model, UMI delegates game state management to the interface (GUI/CLI), while engines focus purely on decision-making. The interface acts as an arbiter, managing:
\begin{itemize}
    \item Tile wall generation and distribution
    \item Turn management and timing
    \item Scoring calculation and application
    \item Game history and replay generation
\end{itemize}

\subsection{Transport Layers}

UMI supports three transport mechanisms:

\subsubsection{Named Pipes (Local)}
Unix named pipes (FIFOs) for local engine communication. Ideal for engine testing and single-machine play.

\subsubsection{TCP Sockets (LAN)}
Direct TCP connections for local network play. Supports client-server and peer-to-peer topologies.

\subsubsection{WebSocket/DHT (Internet)}
Decentralized peer discovery using Kademlia DHT with WebSocket signaling for NAT traversal and global matchmaking.

\section{Protocol Commands}

\subsection{Engine to GUI}

\begin{lstlisting}
id name <x>
id author <x>
id version <x>
umiok
readyok
info <property> <value>
bestmove <move>
error <message>
\end{lstlisting}

\subsection{GUI to Engine}

\begin{lstlisting}
umi
isready
setoption name <id> [value <x>]
position <gamestate>
go [wtime <ms>] [btime <ms>] [byoyomi <ms>]
stop
quit
\end{lstlisting}

\section{Game State Representation}

\subsection{Tile Encoding}

Tiles use a 34-base notation with suit indicators:

\begin{table}[h]
\centering
\begin{tabular}{ll}
\toprule
\textbf{Notation} & \textbf{Meaning} \\
\midrule
m1-m9 & Manzu (Characters) 1-9 \\
p1-p9 & Pinzu (Circles) 1-9 \\
s1-s9 & Souzu (Bamboo) 1-9 \\
E, S, W, N & Winds (East, South, West, North) \\
W, G, R & Dragons (White, Green, Red) \\
\bottomrule
\end{tabular}
\caption{UMI Tile Notation}
\end{table}

\subsection{Position Command Format}

\begin{lstlisting}
position <round> <honba> <kyoutaku> <dealer> <dora> <uradora>
    hands <p1> <p2> <p3> <p4>
    discards <p1> <p2> <p3> <p4>
    calls <p1> <p2> <p3> <p4>
    wall <remaining>
    turn <player>
\end{lstlisting}

Example:
\begin{lstlisting}
position E1 0 0 0 m5
    hands 136m7p89sEENWW 247m25p349sGG
            89m13p67sRRWSS 35m89p12sNNWW
    discards m1,p2,s3,E  ,  ,  
    calls    ,  ,  ,  
    wall 52
    turn 1
\end{lstlisting}

\section{Move Notation}

\subsection{Action Types}

\begin{table}[h]
\centering
\begin{tabular}{lll}
\toprule
\textbf{Move} & \textbf{Format} & \textbf{Description} \\
\midrule
Discard & d<t> & Discard tile t \\
Riichi & r<t> & Declare riichi, discard t \\
Tsumo & t & Self-draw win \\
Ron & ron<p><t> & Win on player p's discard t \\
Pon & pon<p><t><t2> & Call pon on p's t, using t2 from hand \\
Chi & chi<p><t><t1><t2> & Call chi sequence \\
Kan & kan<p><t><t><t><t> & Open/closed kan \\
Pass & p & No action \\
\bottomrule
\end{tabular}
\caption{UMI Move Notation}
\end{table}

\section{DHT-Based Peer Discovery}

\subsection{Kademlia Integration}

UMI uses a modified Kademlia DHT for decentralized peer discovery:

\begin{algorithm}
\caption{UMI Peer Discovery via DHT}
\begin{algorithmic}[1]
\Procedure{JoinNetwork}{bootstrapNodes}
    \State nodeId $\leftarrow$ SHA256(publicKey)
    \State routingTable $\leftarrow$ InitializeKBucket(nodeId)
    \For{each node in bootstrapNodes}
        \State \Call{Ping}{node}
        \State \Call{FindNode}{node, nodeId}
    \EndFor
    \State \Call{RefreshBuckets}{}
\EndProcedure

\Procedure{FindGame}{gameParams}
    \State gameKey $\leftarrow$ \Call{SHA256}{gameParams.rating + gameParams.rules}
    \State closestNodes $\leftarrow$ \Call{Lookup}{gameKey}
    \For{each node in closestNodes}
        \State peers $\leftarrow$ \Call{GetPeers}{node, gameKey}
        \If{peers.found}
            \State \Return \Call{ConnectWebSocket}{peers.address}
        \EndIf
    \EndFor
    \State \Call{Announce}{gameKey, myAddress}
    \State \Return \text{WAITING}
\EndProcedure
\end{algorithmic}
\end{algorithm}

\subsection{WebSocket Signaling}

Once peers are discovered via DHT, WebSocket provides the actual game transport:

\begin{lstlisting}
// Connection establishment
WS /umi/v1/game/{gameId}

// Message format (JSON)
{
    "type": "move",
    "player": 0,
    "move": "d7m",
    "timestamp": 1712345678901,
    "signature": "base64sig"
}
\end{lstlisting}

\section{Security Considerations}

\subsection{Move Verification}

All moves are cryptographically signed using Ed25519. The game state hash chain prevents tampering:

\begin{equation}
H_n = SHA256(H_{n-1} || move_n || timestamp_n || signature_n)
\end{equation}

\subsection{Sybil Resistance}

DHT node IDs require proof-of-work to prevent Sybil attacks:

\begin{equation}
SHA256(nodeId || nonce) < 2^{256 - difficulty}
\end{equation}

\section{Implementation Reference}

\subsection{State Machine}

The UMI engine operates as a finite state machine:

\begin{itemize}
    \item \textbf{IDLE}: Waiting for position command
    \item \textbf{THINKING}: Calculating best move
    \item \textbf{PONDERING}: Opponent's turn, background calculation
    \item \textbf{STOPPED}: Calculation interrupted
\end{itemize}

\section{Conclusion}

UMI provides a robust, extensible protocol for Japanese Mahjong engine communication. By supporting local pipes, TCP sockets, and decentralized DHT/WebSocket networking, it enables use cases ranging from local AI development to global peer-to-peer tournament play without centralized infrastructure.

\begin{thebibliography}{9}
\bibitem{uci} Meyer-Kahlen, S. (2000). Universal Chess Interface. Shredder Computer Chess.
\bibitem{kademlia} Maymounkov, P., \& Mazieres, D. (2002). Kademlia: A peer-to-peer information system based on the XOR metric. IPTPS.
\bibitem{riichi} World Riichi Championship Rules (2022). Standardized Mahjong Rules.
\end{thebibliography}

\end{document}